\section{Transfer to all larger cardinals}
\label{tail:section-transfer}

The common levels of Theorem~\ref{close:common-levels} provide one
cardinal at which all finite geometric shapes have exact boundary
lifting.  This section uses that single cardinal as the base of an
induction through the cardinals.  The levels supplied by the closure
construction themselves need not contain a tail.

The trade between one extra dimension and an increase in cardinality
goes back to Shelah's multidimensional existence and uniqueness methods
\cite[Theorems~5.1--5.2]{Shelah1983B}; compare their later formulation in
\cite[Lemma~11.6]{ShelahVasey2024}. The argument below
is given for the exact boundary maps and cardinal assignments used here.

Fix a common uniform level $\gamma$ from
Theorem~\ref{close:common-levels}, for all the geometric shapes and
their exact boundary assertions.  Choose it strictly above the
no-growth threshold $\lambda$ of
Corollary~\ref{geom:tail-no-growth}, and above all fixed point and
single-model thresholds.  In particular $\gamma>\lambda_0$, the
cardinality of the distinguished model~$M_0$.  In every pointed
model of cardinality at least $\gamma$, model cardinality equals
the cardinality of a basis of its point geometry.

\begin{definition}[Prescribed piece cardinalities]
\label{tail:piece-cardinalities}
Let $L$ be an allowed family on a finite axis set $a$.  A cardinal
assignment on $L$ is a family of cardinals
\[
 (\kappa_s:s\in L),\qquad
 \gamma\leq\kappa_s\leq\kappa_t\quad(s\subseteq t\text{ in }L).
\]
A system with this assignment satisfies the span and inclusion
conditions of Definition~\ref{sys:shapes}, with a witnessing basis
partition $(H_s:s\in L)$ such that $|H_s|=\kappa_s$.
The bases are not named.  A prescribed boundary isomorphism has
the meaning of Definition~\ref{exact:boundary-data}.
\end{definition}

For $t\subseteq a$ let $a_L(t)$ be its largest active subset.
The span identities and monotonicity imply
\begin{equation}
 \rk P(X_t)=\max\{\kappa_s:s\in L,\ s\subseteq t\}
           =\kappa_{a_L(t)}.
 \label{tail:vertex-size}
\end{equation}
All these ranks are at least $\gamma$, so
Corollary~\ref{geom:tail-no-growth} shows that the same formula
gives $\|X_t\|$.  It also shows that an inactive vertex equals
$X_{a_L(t)}$.  Lemma~\ref{sys:relative-bases} applies without a
uniform-cardinality requirement: its proof only uses the finite
shape and the basis axioms of the pregeometry.

\begin{lemma}[Choosing subcubes with prescribed supports]
\label{tail:support-subcubes}
Suppose two systems $\mathcal X^-$ and $\mathcal X^+$ have the
same cardinal assignment, and prescribed isomorphisms $h_t$ on
their proper faces.  Suppose their top is active and put
\[
 \kappa=\max_{s\in L}\kappa_s,\qquad
 Q=\{s\in L:\kappa_s=\kappa\}.
\]
One can choose witnessing global basis partitions
$(H_s^-)$ and $(H_s^+)$ such that
\begin{equation}
 h_s[H_s^-]=H_s^+\qquad(s\in L,\ s\subsetneq a).
 \label{tail:proper-basis-choice}
\end{equation}
Fix such partitions.  Let $\chi$ be a cardinal with
$\gamma\leq\chi<\kappa$ which bounds all $\kappa_s<\kappa$.
Given at most $\chi$ requested elements in the finitely many
vertices on both sides, there are subsystems $\mathcal U^\pm$
inside $\mathcal X^\pm$ with the following properties:
\begin{enumerate}
\item Each requested element belongs to its indicated vertex.
      Every vertex whose size is strictly below $\kappa$ is
      retained in full.
\item Their witnessing pieces are subsets
      $I_s^\pm\subseteq H_s^\pm$.  They equal $H_s^\pm$ when
      $s\notin Q$, and have cardinality $\chi$ when $s\in Q$.
\item On every proper face,
      $h_t[U_t^-]=U_t^+$.
\end{enumerate}
If compatible subsystems have already been chosen whose witnessing
pieces are subsets of these fixed partitions $(H_s^\pm)$, and whose
total vertex sizes are at most $\chi$, their entire vertices can be
included among the requests.  Moreover, at every $s\in Q$ one can
require exactly $\chi$ additional basis coordinates beyond those
old selected pieces, on both sides.  All resulting inclusions
preserve the relative independence condition for system arrows.
\end{lemma}

\begin{proof}
Choose the pieces on the minus side.  For each proper active $s$,
the isomorphism $h_s$ takes its lower span onto the corresponding
lower span on the plus side, since it agrees with the prescribed
maps on every lower vertex.  Thus $h_s[H_s^-]$ is a relative basis
at~$s$.  By Lemma~\ref{sys:relative-bases}, these relative bases
assemble in the plus system.  Complete them by a relative basis
at its top.  Their cardinalities have the required values because
relative rank is invariant under isomorphism.  This gives the
claimed partitions.

We construct the desired subsystems by a countable closure
process, on both sides and at all the finitely many vertices.
Begin with the requested sets, the fixed point, and any specified
old subsystems.  At every growing active label $s\in Q$, also
request $\chi$ basis coordinates.  In the extension version
choose these outside the old selected piece.  There are enough
coordinates since $|H_s^\pm|=\kappa$ and the old piece has size
at most $\chi<\kappa$.  At a proper label the choices may first
be made on the minus side and transported by $h_s$; these are
still fresh on the plus side.  At the top make the choices on
both sides separately.  Retain all nongrowing vertices in full.

At each finite stage choose, by the L\"owenheim--Skolem axiom,
strong submodels of the ambient vertices containing all requests
so far.  Their sizes are at most $\chi$: the LS bound, point,
retained smaller vertices, and requested data all fit within that
cardinal.  At the next stage make the following further requests.
\begin{enumerate}
\item For each inclusion of labels $s\subseteq t$, put all
      elements of the currently chosen model at $s$ into the
      requests at~$t$.
\item For every point obtained in a chosen vertex, choose a finite
      support for it from the union of the global basis pieces
      below that vertex.  Request each supporting coordinate in
      the vertex of its own active label.
\item At each proper face, request the image under $h_t$ of every
      element currently included on the minus side, and the
      inverse image of every element currently included on the
      plus side.
\end{enumerate}
Also retain all previous requests and chosen models.  These
operations make at most $\chi$ requests at each stage.  Each
ambient point set has the given basis partition, so the supports
in the second operation exist and lie in the indicated vertices.
For a retained smaller vertex all supporting labels are retained
as well: if $s\subseteq t$ and $\kappa_t<\kappa$, monotonicity
gives $\kappa_s<\kappa$.  Thus no operation requires enlarging a
retained vertex beyond its given ambient model.

Take the unions of the chosen models through the countable
process.  Coherence makes all successive and vertex inclusions
strong, and the AEC chain axioms give strong limit models
$U_t^\pm\leq_\K X_t^\pm$.  On a proper face the image and
inverse-image requests give exactly $h_t[U_t^-]=U_t^+$.
This construction used the maps one face at a time; it did not
assume that their union on possible extra intersections was
already a function.

Put $I_s^\pm=H_s^\pm\cap P(U_s^\pm)$.  Every point of
$P(U_t^\pm)$ occurs in a finite-stage model, and its supporting
coordinates were subsequently placed at their own labels.
Conversely, the point set of $U_t^\pm$ is relatively closed in
the ambient model and contains all selected coordinates below
$t$.  Therefore
\[
 P(U_t^\pm)=\cl\left(\bigcup_{s\in L,\ s\subseteq t}
                                      I_s^\pm\right).
\]
The selected pieces are subsets of the single global independent
basis and hence are independent together.  They supply precisely
the required decomposition.  Nongrowing pieces and vertices were
retained in full.  At a growing label at least $\chi$ coordinates
were requested and its model has size at most $\chi$, so its
piece has size exactly $\chi$.  In the extension version, at
least $\chi$ fresh coordinates were requested; at most $\chi$
coordinates altogether were added.  The additional piece
therefore has cardinality exactly~$\chi$.

Finally, the selected subsets increase when old subsystems are
included.  Being subsets of the same global basis proves the
relative independence requirement for their inclusions, using
the criterion of Lemma~\ref{sys:relative-bases}.
\end{proof}

The subcube lemma makes compatible smaller systems available, but an
arbitrary isomorphism between the next two subcubes need not extend the
previous one.  The cardinal induction resolves this by adding an axis:
the old full isomorphism becomes part of a boundary prescription at a
strictly smaller maximum cardinal.  This is why the induction must cover
all finite shapes simultaneously.

\begin{theorem}[Exact lifting at every larger cardinal]
\label{tail:all-cardinals}
For every finite allowed shape and every monotone assignment
$(\kappa_s:s\in L)$ of cardinals at least $\gamma$, any prescribed
boundary isomorphism between two systems with that assignment
extends to an isomorphism of their tops.
\end{theorem}

\begin{proof}
We induct on the cardinal
$\kappa=\max_{s\in L}\kappa_s$, proving the assertion
simultaneously for all finite shapes and all assignments with
that maximum.  For $\kappa=\gamma$ all pieces have size
$\gamma$, so Theorem~\ref{exact:geometric-boundary} applies at
the chosen common level.  Suppose $\kappa>\gamma$ and the
assertion is known at every smaller maximum.

As before, an inactive top equals a proper active vertex and its
prescribed map already gives the conclusion.  Assume the top is
active.  Keep all smaller pieces and their vertices in full.
There are only finitely many of them, so their cardinalities
have a common bound below~$\kappa$.  Choose the global basis
partitions of Lemma~\ref{tail:support-subcubes}.

Let $\delta=\cf(\kappa)$.  Choose an increasing sequence
$(\alpha_i:i<\delta)$ of ordinals below $\kappa$, cofinal in
$\kappa$, and fix enumerations of every growing vertex on each
side by~$\kappa$.  We construct continuous chains of subsystems
\[
 (\mathcal U_i^-:i<\delta),\qquad
 (\mathcal U_i^+:i<\delta),
\]
and compatible full system isomorphisms
$f_i:\mathcal U_i^-\simeq\mathcal U_i^+$.
Every proper-face restriction of $f_i$ is the given $h_t$ on
that subface.  Their pieces are increasing subsets of the
chosen global partitions.  Nongrowing vertices are retained
in full; all growing pieces at stage $i$ have the same
cardinality $\rho_i<\kappa$, at least~$\gamma$.

For the initial stage, choose $\chi<\kappa$ bounding $\gamma$,
the retained pieces, and the sizes of the requested initial
enumeration segments.  Apply Lemma~\ref{tail:support-subcubes}
on both sides.  This gives systems with the same piece sizes
and maximum below $\kappa$, whose proper-face maps are the
restrictions of the given prescription.  The induction
hypothesis at that maximum supplies~$f_0$.

Suppose $\mathcal U_i^\pm$ and $f_i$ are given.  Choose a
cardinal $\chi<\kappa$ which bounds $\gamma$, all retained
piece sizes, all old vertex sizes, and the next requested
enumeration segments.  The size checks justifying this choice
at every proper stage are given below.  Apply the extension
part of Lemma~\ref{tail:support-subcubes}.  It produces
$\mathcal U_{i+1}^\pm$ containing the old systems and the
requested segments, with exactly $\chi$ new independent
coordinates at every label in
\[
 Q=\{s\in L:\kappa_s=\kappa\},
\]
and none at the other labels.  The set $Q$ is upward closed:
if $s\in Q$ and $s\subseteq t\in L$, then
$\kappa\leq\kappa_t\leq\kappa$, so $t\in Q$.

Add a new last axis $k$.  Regard
$\mathcal U_i^\pm\to\mathcal U_{i+1}^\pm$ as the two layers
of systems with active labels
\[
 L'=L\cup\{s\cup\{k\}:s\in Q\}.
\]
The old pieces have their previous cardinalities and the
additional pieces all have cardinality~$\chi$.  These are
valid decompositions: the old and added coordinate sets are
disjoint subsets of the fixed global basis, and their unions
span the new vertices.  The cardinal assignment on $L'$ is
monotone.  Comparisons among old labels use the old
monotonicity, comparisons among new labels use equality to
$\chi$, and an old label below a new one has size at most
$\chi$ by its choice.  No new label is contained in an old
one.  All assigned cardinals are at least $\gamma$ and the
maximum is still strictly below~$\kappa$.

Prescribe the boundary maps of these enlarged shapes by $f_i$
on the entire old cube and by the given $h_t$ on every enlarged
proper face of the new cube.  They commute on all indicated
subfaces: on an old proper face the two prescriptions already
agree by the invariant for~$f_i$.  The cardinal induction
hypothesis applies to this shape as well, since it is
simultaneous in the number of axes and its maximum decreased.
It yields an isomorphism of the new tops extending $f_i$ and
all the enlarged proper-face maps.  This is the required
$f_{i+1}$.

At a nonzero limit $i<\delta$, take componentwise unions and
put $f_i=\bigcup_{j<i}f_j$.  The maps are coherent, so $f_i$
is an onto isomorphism between these unions, with its required
face restrictions.  The component unions are strong submodels
of the actual ambient vertices by the AEC common-upper-bound
axiom.  Union pieces are still independent by finite
character.  Every point in a union vertex occurs at some
stage and has a finite support in that stage's selected
pieces, so the union pieces span.  These are thus systems
of the required shape, and inclusions continue to preserve
relative independence.

Here are the cardinal calculations at proper stages.  If
$\kappa=\tau^+$ is a successor cardinal, choose
$\chi=\tau$ at the initial and every growth stage.  Every
proper stage index has cardinality at most $\tau$, and every
union at such a stage has size at most
$\tau\cdot\tau=\tau$.  The enumeration segments below
$\kappa$ also have size at most $\tau$.  Hence every stated
choice and every proper-stage induction application has
maximum strictly below~$\kappa$.

If $\kappa$ is a regular limit cardinal, there are fewer than
$\kappa$ earlier stages at every $i<\delta=\kappa$.
Regularity bounds their sizes and their union by a cardinal
below~$\kappa$.  Enlarge the next $\chi$ to cover that bound,
the finitely many retained sizes, and the next enumeration
segments.  It can still be chosen below the limit cardinal
$\kappa$, and may be increased cofinally through~$\kappa$.

If $\kappa$ is singular, at a proper stage $i<\delta$ fewer
than $\cf(\kappa)$ earlier cardinals have been used.  Their
supremum is therefore below~$\kappa$.  Both that supremum
and $|i|$ have a common cardinal bound below~$\kappa$, as do
the finitely many retained sizes and the next enumeration
segments.  Their maximum bounds the union and can be used
in choosing~$\chi$.  In particular this argument applies
when $\cf(\kappa)=\aleph_0$; the final countable union has
size $\kappa$, while every proper stage remains below it.

The variable pieces have equal sizes on the two sides and
at all growing labels.  This is immediate initially and at
successors, where the same additional cardinal $\chi$ is
used everywhere.  At a limit their sizes are the same
cardinal sum of the initial piece and the previously added
disjoint pieces.  The preceding bounds keep that sum below
$\kappa$ at a proper stage.  Monotonicity also persists,
since the smaller retained pieces were bounded by the
initial variable size.

Finally take the unions over $i<\delta$.  The requested
enumeration segments are cofinal, so these unions are the
two original systems.  The union of the $f_i$ is an onto
isomorphism of the tops with every prescribed restriction.
The final union is allowed to have cardinality $\kappa$;
the induction hypothesis was used only at proper stages,
where the maximum was smaller.  This completes the cardinal
induction.
\end{proof}

\begin{corollary}[Pointed categoricity on a tail]
\label{tail:pointed-categoricity}
For every cardinal $\kappa\geq\gamma$, any two models of
cardinality $\kappa$ carrying the fixed point pattern are
pointedly isomorphic.
\end{corollary}

\begin{proof}
On zero axes the only active piece is a basis of the whole
point geometry.  Its size is the model cardinality by
Corollary~\ref{geom:tail-no-growth}.  Apply
Theorem~\ref{tail:all-cardinals} with the one-element
cardinal assignment~$\kappa$.  The boundary is empty.
\end{proof}

\begin{theorem}[Categoricity on a tail]
\label{tail:categoricity}
If an AEC $\K$ is categorical in unboundedly many cardinals,
then there is a cardinal $\theta$ such that $\K$ is
categorical in every cardinal $\kappa\geq\theta$.
Here categoricity includes the existence of a model at
each such cardinal.
\end{theorem}

\begin{proof}
All preceding constructions were made for this fixed AEC
under the hypothesis of unbounded categoricity.  Choose
$\gamma$ as above.  Fix a sufficiently large categorical
cardinal $\rho\geq\gamma$ at which a pointed presentation
with the chosen diagram of $M_0$ exists and the stabilized
$\K$-diagrams on its enumeration are preserved by strong
maps into larger models.  Such a $\rho$ is available from
the single-model working levels and
Proposition~\ref{geom:pointed-levels}.  We may take
$\theta=\rho$.

Let $M$ be any model of cardinality
$\kappa\geq\rho$.  Choose a subset of its universe of
cardinality~$\rho$.  The LS axiom gives a strong submodel
$N\leq_\K M$ containing that subset and of size at most
$\rho+\LS(\K)=\rho$.  Thus $\|N\|=\rho$ exactly.
Categoricity in $\rho$ identifies $N$ with the pointed
presentation just chosen.  Transport its copy of $M_0$
into $N$ and then into~$M$.  The fixed diagram on its
whole enumeration is preserved by the strong inclusion,
so $M$ has the required point pattern.  Notice that this
uses categoricity at the fixed smaller cardinal $\rho$,
not a categoricity assertion at~$\kappa$.

Consequently any two models of cardinality
$\kappa\geq\theta$ can be given the chosen point pattern.
Corollary~\ref{tail:pointed-categoricity} identifies them
pointedly, and forgetting that point gives a
$\K$-isomorphism.  This proves uniqueness.

For existence, take a linear order $I$ of cardinality
$\kappa$.  The order presentation of
Theorem~\ref{found:order-template} is a member of $\K$ and
has size exactly~$\kappa$: its distinct skeleton gives
the lower bound, and its at-most-$\LS(\K)$ labels on
finite lists give the upper bound
$\kappa+\LS(\K)=\kappa$.  Hence there is exactly one
isomorphism type, rather than merely at most one, at every
cardinal in the asserted tail.
\end{proof}

\subsection*{Conclusion of the uniform theorem}
\label{tail:uniform-threshold}

\begin{proof}[Proof of Theorems~\ref{main:qualitative} and
  \ref{main:uniform}]
Theorem~\ref{tail:categoricity} proves the qualitative
assertion of Theorem~\ref{main:qualitative}.  Apply the
set-coding reduction of
Proposition~\ref{found:uniform-reduction} at a fixed infinite
LS bound~$\lambda$.  For each of its set-many AEC codes,
either the categoricity spectrum is bounded, or the
qualitative assertion gives categoricity on a tail.
In the bounded case choose a strict upper bound for every
categorical cardinal of that code.  In the unbounded case
choose a threshold for its categorical tail.  Replacement
over the set of codes gives a common cardinal bound
$\mu(\lambda)\geq\lambda$ above all these choices.

If an AEC with LS number at most $\lambda$ is categorical
at any cardinal $\kappa\geq\mu(\lambda)$, its code cannot
be one of the bounded-spectrum cases.  Its categorical
tail therefore contains every
$\kappa'\geq\mu(\lambda)$.  This is exactly
Theorem~\ref{main:uniform}, with the same threshold in
its premise and conclusion.  No regularity, successor,
amalgamation, or no-maximal-model hypothesis has been
added to the statement.
\end{proof}
